\chapter{Коли записувати, а коли читати: додаємо \nt{ACET}}
\chaptermark{Додаємо \nt{ACET}}
\label{sec:acet}

\section{Чого бракує}

Мікросхема пам'яті, крім проводів адреси і даних, має ще одну лінію --- дозвіл запису. Поки вона вимкнена, пам'ять лише читається; коли ввімкнена, те, що є на лініях даних, записується. Без такої лінії пам'ять не працює: якщо кожне читання водночас щось записує, прочитане, зокрема прочитане з помилкою, одразу ж записується назад.

У мозку ця задача постає гостріше, ніж у мікросхемі. У розділі~\ref{sec:glutcomputer} пам'яттю була група \nt{GLUT}-нейронів, яку тримають увімкненою її власні зв'язки (рис.~\ref{fig:bistable}). У розділі~\ref{sec:dofa} ми бачили, що ці зв'язки навчаються правилом Гебба просто під час роботи. Отже, ті самі синапси і зберігають старе, і приймають нове, а окремої фази запису, як і годинника, немає. Чим мережа відрізняє момент, коли треба запам'ятати те, що бачиш, від моменту, коли треба згадати те, що знаєш? У розділі~\ref{sec:neurontypes} ми припустили, що цю лінію тримає \nt{ACET}. У цьому розділі ми перевіримо, як така лінія має працювати і чому без неї не обійтися.

\section{Три дії одного проводу}

Ацетилхолінових нейронів, що працюють усередині мозку, небагато, і зібрані вони в кількох невеликих групах, а їхні виходи розходяться по всій корі. Це широкомовний канал, такий самий, як \nt{DOPA} і \nt{NORA}. Рівень ацетилхоліну в корі високий в уважному неспанні і низький у повільному сні~\cite{Hasselmo1999}. Як і в дофаміну, зміст сигналу задає рецептор отримувача (твердження~\ref{prop:receptor}): ацетилхолін має і швидкі рецептори, що відкривають канал, і повільні, що запускають реакції всередині клітини. Для нас важливі три дії.

\emph{Перша: послабити внутрішні зв'язки.} Досліди на зрізах нюхової кори показали, що ацетилхолін сильно пригнічує передачу зв'язками між нейронами однієї області і майже не зачіпає входи, що надходять у неї ззовні~\cite{HasselmoBower1992}.

\emph{Друга: посилити входи.} Ацетилхолін підвищує збудливість нейронів і полегшує передачу деякими вхідними шляхами; сигнал від органів чуттів стає гучнішим за власну активність мережі~\cite{Hasselmo2006}.

\emph{Третя: полегшити зміну ваг.} За високого рівня ацетилхоліну зв'язки змінюються легше~\cite{Hasselmo2006}.

Для інженера всі три дії --- одна операція: перемикання з режиму «читати» в режим «записувати». Мережа перестає слухати саму себе, починає слухати вхід і запам'ятовує те, що чує.

\section{Мережа, яка дописує}

Візьмімо $N$ \nt{GLUT}-нейронів, з'єднаних один з одним. Збережімо в їхніх зв'язках кілька образів --- наборів по $pN$ одночасно активних нейронів --- правилом Гебба~\cite{Hebb1949}. Якщо подати на таку мережу половину образу, зв'язки збудять відсутню половину: мережа \emph{дописує} образ за частиною, як група на рис.~\ref{fig:bistable} підтримувала сама себе. Це асоціативна пам'ять: адресою слугує частина вмісту. \nt{GABA}, як у розділі~\ref{sec:gaba}, тримає загальну кількість активних нейронів близько $pN$, щоб разом не могли засвітитися два образи.

Позначмо рівень ацетилхоліну $a$, від $0$ до $1$, і запишімо три його дії в модель прямо:
\begin{empheq}[box=\keybox]{equation}
\tau\frac{\dd r_i}{\dd t} = -r_i + F\Bigl(\tfrac{1+a}{2}\,x_i + (1-a)\sum_j W_{ij}\,r_j - g\Bigr),
\qquad
\frac{\dd W_{ij}}{\dd t} = \eta\,a\,(r_i-p)(r_j-p) .
\label{eq:acetnet}
\end{empheq}
Тут $r_i$ --- частота нейрона $i$, $x_i$ --- його вхід від органів чуттів, $F$ --- порогова передавальна характеристика, $g$ --- загальне гальмування від \nt{GABA}, що зростає, коли активних нейронів більше за $pN$. Множник $\frac{1+a}{2}$ --- посилення входу, $1-a$ --- послаблення внутрішніх зв'язків, $\eta a$ --- швидкість навчання. Друге рівняння --- правило Гебба~\eqref{eq:hebb} у формі, зручній для рідкісних образів~\cite{Tsodyks1988}: вага росте, коли обидва нейрони активніші, ніж зазвичай, і зменшується, коли активний лише один. За від'ємною частиною ваг, як завжди, стоять \nt{GABA}-посередники. Годинника немає: і нейрони, і ваги змінюються неперервно.

\begin{figure}[t]
\centering
\begin{tikzpicture}[>=Stealth,
  nrn/.style={circle,draw,very thick,fill=blue!6,minimum size=0.8cm,inner sep=0pt,font=\small},
  acet/.style={circle,draw,very thick,double,double distance=1.2pt,fill=orange!15,minimum size=1.35cm,inner sep=0pt,font=\scriptsize},
  bc/.style={->,thick,densely dashed,orange!85!black}]
\foreach \i/\y in {1/2.4,2/1.2,3/0}{
  \node[nrn] (N\i) at (3,\y) {$r_\i$};
  \draw[->,thick,blue!60!black] (0,\y) node[left,black] {\small $x_\i$} -- (N\i);}
\node[left,align=right,font=\small] at (-0.5,1.2) {від органів\\чуттів};
\draw[->,thick,gray!60!black] (N1) to[bend left=30] (N2);
\draw[->,thick,gray!60!black] (N2) to[bend left=30] (N1);
\draw[->,thick,gray!60!black] (N2) to[bend left=30] (N3);
\draw[->,thick,gray!60!black] (N3) to[bend left=30] (N2);
\draw[->,thick,gray!60!black] (N1.east) .. controls (4.6,2.0) and (4.6,0.4) .. (N3.east);
\node[right,font=\small,gray!40!black,align=left] at (4.7,1.2) {внутрішні\\зв'язки $W$};
\node[acet] (A) at (1.2,-1.9) {\nt{ACET}};
\draw[bc] (A) -- (1.6,-0.15);
\node[font=\small,orange!60!black,anchor=east] at (0.95,-0.9) {$+$ вхід};
\draw[bc] (A) -- (4.3,0.6);
\node[font=\small,orange!60!black,anchor=west] at (4.3,0.1) {$-$ внутрішні зв'язки, $+$ швидкість навчання};
\node[right,font=\small,align=left] at (2.2,-1.9) {високий рівень: «записувати»\\низький рівень: «читати»};
\end{tikzpicture}
\caption{\nt{ACET} як лінія дозволу запису. Нейрони $r_i$ отримують вхід $x_i$ від органів чуттів (сині стрілки) і збуджують один одного через внутрішні зв'язки $W$ (сірі), у яких зберігаються образи. Ацетилхолін (штрихові стрілки) посилює вхід, послаблює внутрішні зв'язки і полегшує зміну ваг, як у~\eqref{eq:acetnet}.}
\label{fig:acetnet}
\end{figure}

\section{Запис і читання заважають одне одному}

Перевірмо модель~\eqref{eq:acetnet} на задачі, в якій запис і читання справді конфліктують. Мережа з $200$ нейронів уже зберігає п'ять образів по $20$ активних нейронів. Тепер треба запам'ятати шостий, який наполовину збігається з одним зі старих: десять нейронів у них спільні. Так буває постійно: нова кімната схожа на знайому, нове обличчя --- на старе.

\emph{Запис за низького $a$.} Спочатку відокремимо перші дві дії ацетилхоліну від третьої: залишимо швидкість навчання повною, а змінюватимуться лише посилення входу і послаблення внутрішніх зв'язків. При $a=0$ вхід запалює спільні десять нейронів, а вони через внутрішні зв'язки запалюють решту половини \emph{старого} образу. \nt{GABA} не дає горіти обом одразу, і старий образ, підтриманий власними зв'язками, витісняє новий, який тримається лише на вході. Мережа бачить не те, що перед нею, а те, що пам'ятає, --- і правило Гебба старанно записує побачене.

У підсумку новий образ узагалі не запам'ятався: за його відмінною половиною мережа нічого не згадує. А старий зіпсований: викликаний своєю відмінною половиною, він тягне за собою в середньому шість чужих нейронів із десяти. При $a=0.1$ новий образ уже запам'ятовується, але зливається зі старим, і псування навіть сильніше: кожен із двох, викликаний своєю половиною, тягне за собою близько восьми чужих нейронів; починаючи з $a=0.2$ запис чистий: зайвих нейронів під час згадування менше за один (усереднення за десятьма прогонами).

\emph{Запис зі справжньою швидкістю навчання.} Тепер нехай ацетилхолін, як у~\eqref{eq:acetnet}, задає і швидкість навчання. За одне пред'явлення новий образ запам'ятовується цілком лише при $a\ge0.9$; при $a=0.75$ --- на вісім десятих, при $a\le0.5$ не запам'ятовується зовсім.

\emph{Читання.} Подамо на мережу з п'ятьма чисто записаними образами половину образу, а потім приберемо і її. При $a\le0.2$ мережа дописує образ цілком і тримає його сама; при $a=0.3$ --- майже цілком; при $a\ge0.4$ внутрішні зв'язки надто послаблені, і після того як підказка зникла, активність гасне (рис.~\ref{fig:acetsim}).

\begin{figure}[t]
\centering
\begin{tikzpicture}[>=Stealth,x=9cm,y=3.2cm]
\draw[->,gray!70!black] (0,0) -- (1.1,0) node[right,black] {\small $a$};
\draw[->,gray!70!black] (0,0) -- (0,1.2);
\foreach \x in {0,0.2,0.4,0.6,0.8,1.0}{\draw[gray!60!black] (\x,-0.02) -- (\x,0.02); \node[below,font=\scriptsize] at (\x,0) {$\x$};}
\foreach \y in {0,0.5,1}{\draw[gray!60!black] (-0.01,\y) -- (0.01,\y); \node[left,font=\scriptsize] at (0,\y) {$\y$};}
\node[rotate=90,font=\small] at (-0.1,0.6) {частка відновленого образу};
\draw[very thick,blue!60!black] plot[mark=*,mark size=1.6pt] coordinates {(0,1.00) (0.1,1.00) (0.2,1.00) (0.3,0.96) (0.4,0.04) (0.5,0.00) (0.75,0.00) (1.0,0.00)};
\draw[very thick,red!70!black] plot[mark=*,mark size=1.6pt] coordinates {(0.1,0.00) (0.2,0.00) (0.3,0.00) (0.5,0.00) (0.75,0.80) (0.9,1.00) (1.0,1.00)};
\node[font=\small,blue!60!black,anchor=west,align=left] at (0.03,0.5) {читання:\\образ\\за половиною};
\node[font=\small,red!70!black,anchor=east,align=right] at (1.0,0.35) {запис:\\новий образ\\за раз};
\end{tikzpicture}
\caption{Модель~\eqref{eq:acetnet}: $200$ нейронів, образи по $20$ активних, середнє за десятьма прогонами. Сині точки --- читання: яку частку образу відновлено за його половиною після того, як підказку прибрано. Червоні точки --- запис: яка частка нового образу, наполовину схожого на старий, згадується після одного пред'явлення, якщо ацетилхолін задає і швидкість навчання. Читання працює при $a\le0.3$, запис --- при $a\ge0.9$.}
\label{fig:acetsim}
\end{figure}

\begin{keyblock}
\begin{proposition}[Запис і читання потребують різних режимів]
\label{prop:writeread}
У моделі~\eqref{eq:acetnet}, де ті самі зв'язки зберігають старе і навчаються нового, а ацетилхолін задає і швидкість навчання, немає рівня $a$, за якого однаково добре йдуть і запис, і читання. Для запису внутрішні зв'язки треба послабити, інакше мережа запише не вхід, а згадане; для читання їх треба посилити, інакше вона не допише образ за частиною. Потрібен перемикач, і ним може слугувати один широкомовний провід.
\end{proposition}
\end{keyblock}

Якби ацетилхолін не змінював швидкості навчання, для обох справ годилася б вузька смуга $0.2\le a\le0.3$. Але тоді мережа навчалася б і під час кожного читання, тобто заново записувала б прочитане разом із його помилками. Саму ідею --- пригнічувати внутрішні зв'язки під час навчання --- взято з моделей, побудованих за вимірюваннями на зрізах~\cite{HasselmoAndersonBower1992}. Числа вище стосуються нашої спрощеної моделі; де саме лежать межі режимів у мозку, невідомо.

\section{Наскільки вірити очам}

Перемикач «записувати/читати» --- крайній випадок загальнішого питання: наскільки вірити входу і наскільки --- пам'яті? На це питання інженери мають точну відповідь, і вона пояснює, чому один провід керує водночас і режимом, і швидкістю навчання.

Нехай мережа стежить за величиною $s(t)$, яка повільно блукає: за секунду її дисперсія зростає на $q$. Органи чуттів повідомляють $y(t)=s(t)+$шум з інтенсивністю $R$. Найкращу оцінку $\hat s$ у неперервному часі дає фільтр Калмана--Б'юсі~\cite{KalmanBucy1961}:
\begin{empheq}[box=\keybox]{equation}
\frac{\dd\hat s}{\dd t} \;=\; \frac{P}{R}\,\bigl(y-\hat s\bigr),
\qquad
\frac{\dd P}{\dd t} \;=\; q-\frac{P^2}{R} ,
\label{eq:kalman}
\end{empheq}
де $P$ --- дисперсія помилки оцінки, тобто наскільки мережа сама не впевнена в тому, що знає. Перше рівняння --- те саме RC-коло, що й у мембрани в моделі нейрона~\eqref{eq:glutneuron}, але зі сталою часу $R/P$, яка змінюється. Коли мережа не впевнена, $P$ велике, стала часу мала, і оцінка швидко йде за входом: це «записувати». Коли впевнена, $P$ мале, і оцінка майже не слухає входу, тримаючись за пам'ять: це «читати».

В усталеному режимі $P=\sqrt{qR}$ і стала часу пам'яті дорівнює $\sqrt{R/q}$: якщо світ відходить на одиницю за сто секунд, $q=0.01$, а шум органів чуттів $R=1$, пам'ять має пам'ятати близько десяти секунд.

Головне тут те, що одне число $P/R$ задає водночас дві речі: і вагу входу порівняно з пам'яттю, і швидкість, з якою змінюється збережена оцінка. Це саме те, що робить ацетилхолін у~\eqref{eq:acetnet}. За гіпотезою~\cite{YuDayan2005}, ацетилхолін і повідомляє мережі величину, подібну до $P$, --- \emph{очікувану} невизначеність, ту, про яку мережа знає заздалегідь. Повільним цей канал буває не завжди: частина \nt{ACET}-нейронів відповідає на винагороду і покарання за два десятки мілісекунд, тим сильніше, чим несподіваніше підкріплення~\cite{Hangya2015}.

\begin{keyblock}
\begin{proposition}[Одна ручка для ваги входу і швидкості навчання]
\label{prop:kalman}
В оптимальному фільтрі~\eqref{eq:kalman} вага, з якою вхід змішується з пам'яттю, і швидкість навчання --- та сама величина $P/R$. Тому розумно керувати ними одним сигналом. За незмінних властивостей світу цей сигнал виходить на сталий рівень $\sqrt{q/R}$, і підлаштовувати його треба, лише коли властивості світу змінюються.
\end{proposition}
\end{keyblock}

Друге речення твердження пояснює результат розділу~\ref{sec:nora}, де підлаштування швидкості навчання за несподіваністю не виграло в найкращої сталої: у світі з незмінними властивостями змін найкраща швидкість навчання і є сталою. Виграш з'являється, коли світ раптом стає іншим. Тоді оцінка несподівано опиняється далеко від входу, а $P$ про це не знає. Правильна дія --- різко підняти $P$ і тим самим перейти в режим запису. За гіпотезою, яку ми обговорювали в розділі~\ref{sec:nora}, про таку \emph{неочікувану} зміну повідомляє \nt{NORA}~\cite{YuDayan2005}. Розподіл праці виходить природним: \nt{NORA} помічає, що світ змінився, \nt{ACET} тримає мережу в режимі запису, доки нову обстановку не вивчено.

\section{Сон як режим читання}

Якщо високий рівень ацетилхоліну --- режим запису, то що робить мережа за низького? У повільному сні рівень ацетилхоліну падає, внутрішні зв'язки звільняються від пригнічення, входу від органів чуттів майже немає. Мережа в режимі читання без підказки програє те, що в ній записано. Реєстрація активності показує, що нейрони, які вдень працювали разом, у повільному сні знову спрацьовують разом~\cite{WilsonMcNaughton1994}, і навіть у тому самому порядку~\cite{Skaggs1996}. За гіпотезою~\cite{Hasselmo1999}, ці повтори переписують швидко вивчене вдень у довготривалу пам'ять (штучне подовження коротких спалахів повторів покращує пам'ять~\cite{FernandezRuiz2019}): удень --- запис із входу, уночі --- перезапис зсередини. Для інженера це схоже на запис у швидкий буфер удень і його скидання в повільну постійну пам'ять уночі.

\section{Підсумок}

\begin{table}[t]
\centering
\footnotesize
\setlength{\tabcolsep}{4pt}
\renewcommand{\arraystretch}{1.15}
\begin{tabular}{>{\raggedright}p{3.2cm}>{\raggedright}p{4.2cm}>{\raggedright\arraybackslash}p{6.0cm}}
\toprule
Що робить \nt{ACET} & Як & Що відомо \\
\midrule
послаблює внутрішні зв'язки & множник $1-a$ у~\eqref{eq:acetnet} & виміряно на зрізах \\
посилює вхід & множник $\frac{1+a}{2}$ & підвищення збудливості і передачі вхідними шляхами виміряно \\
полегшує навчання & швидкість $\eta a$ & виміряно \\
перемикає «записувати/читати» & твердження~\ref{prop:writeread} & випливає з моделей; прямого вимірювання режимів немає \\
повідомляє очікувану невизначеність & $P$ у~\eqref{eq:kalman} & гіпотеза \\
звільняє мережу для повторів у сні & низьке $a$ в повільному сні & низький рівень уві сні і повтори виміряно; перезапис у довготривалу пам'ять --- гіпотеза \\
\bottomrule
\end{tabular}
\caption{Що додає \nt{ACET} до мережі з \nt{GLUT}, \nt{GABA}, \nt{DOPA} і \nt{NORA}.}
\label{tab:acet}
\end{table}

\nt{DOPA} каже мережі, куди змінювати ваги, \nt{NORA} --- наскільки рішуче користуватися тим, що вона знає, а \nt{ACET} --- чи слухати їй світ, чи себе (таблиця~\ref{tab:acet}). Це остання широкомовна лінія керування з таблиці~\ref{tab:types}: дозвіл запису, без якого пам'ять, що зберігає образи в тих самих зв'язках, якими вона їх читає, псувала б себе під час кожного читання.

\section{Задачі}

\begin{exercise}[\lvlA]
\label{ex:09:1}
\emph{Скільки пам'ятати.} Фільтр~\eqref{eq:kalman} з $q=0.01$, $R=1$ пам'ятає близько $10$~с. Знайдіть $P_\infty$ і пам'ять $\sqrt{R/q}$, якщо (а)~шум датчика зріс удвічі за амплітудою; (б)~світ став блукати в сто разів швидше. (в)~У скільки разів має зрости шум, щоб мережа пам'ятала хвилину?
\end{exercise}

\begin{exercise}[\lvlB]
\label{ex:09:2}
\emph{Режим запису після зміни.} Світ змінився, і невизначеність фільтра~\eqref{eq:kalman} з $q=0.01$, $R=1$ підскочила до $P_0\to\infty$. (а)~З якою сталою часу $P$ повертається до $P_\infty$? (б)~Через скільки секунд швидкість навчання перевищуватиме усталену не більше ніж на $10\,\%$?
\end{exercise}

\begin{exercise}[\lvlB]
\label{ex:09:3}
\emph{Стала швидкість навчання.} Мережа навчається як $\dd\hat s/\dd t=(y-\hat s)/T$; світ блукає, $\dd s/\dd t=w$, $y=s+n$, де $w$ і $n$ --- білі шуми з інтенсивностями $q$ і $R$. Знайдіть найкраще $T$ і дисперсію помилки за нього. У скільки разів вона зросте, якщо $T$ помилкове в $2$ і в $10$ разів?
\end{exercise}

\begin{exercise}[\lvlB]
\label{ex:09:4}
\emph{Де межа запису.} У \texttt{models/\allowbreak acet\_memory.py} поріг $\theta=0.35$, вага всередині образу $w=0.041$ (в ідеалі $0.045$). Новий образ ділить $m=10$ нейронів зі старим. За якого $a$ у~\eqref{eq:acetnet} решта нейронів старого образу не запалюється від цих $m$ (поріг різкий)?
\end{exercise}

\begin{exercise}[\lvlB]
\label{ex:09:5}
\emph{Моделювання: ємність пам'яті.} Змініть у \texttt{models/\allowbreak acet\_memory.py} функцію \texttt{read\_sweep}: при $a=1$ запишіть $M$ образів ($M$ від $5$ до $100$), потім при $a=0$ прочитайте перші десять за половиною. За якого $M$ з'являються помилки і як граничне $M_{\max}$ залежить від $N$ і $p$?
\end{exercise}

\begin{exercise}[\lvlB]
\label{ex:09:6}
\emph{Проєктування: запис без витоку.} Зі швидкістю навчання $\eta a$ мережа навчається і під час читання. Запропонуйте монотонну $\eta(a)\le\eta$, що дорівнює нулю при $a\le0.3$ і не гірша за $\eta a$ при $a\ge0.9$. Перевірте: після $20$ читань при $a=0.2$ з підказкою, де три чужі нейрони, скільки їх увійшло в образ?
\end{exercise}

\begin{exercise}[\lvlC]
\label{ex:09:7}
\emph{Як переписати буфер уночі.} (а)~Уві сні $a=0.05$, повтор триває $0.1\,\text{с}$ проти $0.2\,\text{с}$ удень при $a=1$. Скільки повторів накопичать зміну ваг одного денного пред'явлення за швидкості $\eta a$ і за $\eta(a)$ задачі~\ref{ex:09:6}? (б)~Сховище навчається в $100$ разів повільніше за буфер і чує образ $20$ разів за ніч по $0.1\,\text{с}$. Скільки потрібно ночей?
\end{exercise}
